\documentclass[11pt,a4paper]{article}
\usepackage[utf8]{inputenc}
\usepackage[margin=1in]{geometry}
\usepackage{amsmath,amssymb}
\usepackage{booktabs}
\usepackage{xcolor}
\usepackage{enumitem}
\usepackage{tcolorbox}
\usepackage{fancyhdr}
\usepackage{hyperref}

\hypersetup{
    colorlinks=true,
    linkcolor=blue!80!black,
    urlcolor=blue!80!black,
    citecolor=blue!80!black
}

\tcbuselibrary{skins,breakable}

% Question box
\newtcolorbox{questionbox}[1][]{
  colback=gray!5!white,
  colframe=gray!60!black,
  fonttitle=\bfseries\small,
  title={Judge Question},
  #1
}

% Answer box
\newtcolorbox{answerbox}[1][]{
  colback=blue!5!white,
  colframe=blue!75!black,
  fonttitle=\bfseries\small,
  title={Your Answer},
  #1
}

% Tip box (amber)
\newtcolorbox{tipbox}[1][]{
  colback=orange!5!white,
  colframe=orange!70!black,
  fonttitle=\bfseries\small,
  title={Tip for Presenting This},
  #1
}

\pagestyle{fancy}
\fancyhf{}
\rhead{\textbf{Space-Guard} $\cdot$ Judge Questions Guide}
\lhead{Smart India Hackathon 2026}
\rfoot{Page \thepage}

\title{\Huge \textbf{Space-Guard: Judge Q\&A Cheat Sheet}\\[0.4em]
\Large \textbf{16 Questions Judges Will Ask --- And How to Answer Them}\\[0.5em]
\large \textbf{Smart India Hackathon 2026 --- Software Edition}}
\author{\textbf{Team Space-Guard}}
\date{\today}

\begin{document}

\maketitle
\thispagestyle{empty}

\begin{tcolorbox}[colback=green!5!white,colframe=green!60!black,title=\textbf{How to Use This Guide}]
This document has 16 questions that judges are very likely to ask you. Each question has:
\begin{itemize}
    \item A simple, clear answer you can say out loud.
    \item A tip on how to say it confidently.
\end{itemize}
Read through this the night before. You do not need to memorize exact words --- just understand the idea, and say it in your own natural way. If you do not know an answer, it is okay to say: ``That is a great point. In the current version we handle it this way\ldots and in Phase 2 we plan to address that fully.''
\end{tcolorbox}

\vspace{1em}
\tableofcontents
\newpage

%% ============================================================
\section{Questions About the Problem}
%% ============================================================

\subsection*{Q1: Is this problem real? Have satellites actually collided?}

\begin{questionbox}
``Has there actually been a real satellite collision? Or is this just a theoretical problem you invented?''
\end{questionbox}

\begin{answerbox}
Yes, absolutely. The most famous example is the \textbf{2009 Iridium--Cosmos collision}. On 10 February 2009, an active American phone satellite called Iridium 33 crashed directly into a dead Russian military satellite called Cosmos 2251, at 789 km above Siberia.

They hit each other at over 50,000 km/h. Both satellites were completely destroyed in milliseconds. The crash created over 2,000 tracked pieces of debris that are still in orbit today, threatening the International Space Station and every other satellite at that height.

This was the first-ever accidental collision between two satellites. And no system warned operators in time to prevent it. Space-Guard is designed to be that warning system.
\end{answerbox}

\begin{tipbox}
This is your strongest answer --- deliver it with confidence. You know this story very well. Make eye contact, slow down your speech, and let the facts land. The numbers are very striking: 50,000 km/h, both satellites destroyed, debris still there today. Say them clearly.
\end{tipbox}

\newpage
\subsection*{Q2: Why doesn't ISRO or NASA already have this kind of system?}

\begin{questionbox}
``If this problem is so serious, why don't existing space agencies already solve it?''
\end{questionbox}

\begin{answerbox}
They do have systems --- but they are slow, expensive, and manual. NASA has the ``CARA'' team. ISRO has ``IS4OM.'' But these systems have limitations:
\begin{itemize}
    \item They take minutes to hours to process the full catalog.
    \item They produce thousands of warnings per day, most of them false alarms, which operators start to ignore.
    \item When a real threat is found, engineers still calculate the escape plan manually, which takes more time.
\end{itemize}
Space-Guard is not replacing these organizations. It is a tool that can plug into their pipelines to make the screening step faster, the risk number more precise, and the escape plan automatic --- all in under 5 seconds.
\end{answerbox}

\begin{tipbox}
Do not say ``NASA does not have this.'' That sounds arrogant and is not true. Say we are \textit{faster, smarter, and automatic} compared to current tools. Show respect for existing work, while clearly explaining what gap we fill.
\end{tipbox}

\newpage
%% ============================================================
\section{Questions About Our Solution}
%% ============================================================

\subsection*{Q3: How accurate is your system? Can we trust it?}

\begin{questionbox}
``How do you know your calculations are correct? Did you test this against real data?''
\end{questionbox}

\begin{answerbox}
Yes. We tested our system against the actual 2009 Iridium--Cosmos collision using real satellite tracking data from that time period.

Without telling the system what to look for, it automatically identified the close pass, calculated a collision time of 16:56 UTC on 10 February 2009 (which exactly matches the historical event), calculated the collision happened at 789 km over Siberia (exact match), and assigned it a ``Critical'' risk rating 48 hours in advance.

This is our proof that the physics and math in our system are correct. We are not making anything up --- we matched the real event.
\end{answerbox}

\begin{tipbox}
This is your killer proof point. Say: ``We did not just test on fake data. We ran our system on the actual 2009 satellite data and it correctly identified the crash that really happened.'' If you have the demo open, click to Step 3 and show the collision animation to back this up visually.
\end{tipbox}

\vspace{1em}
\subsection*{Q4: How does your system decide when a situation is dangerous?}

\begin{questionbox}
``What is your definition of a dangerous conjunction? How do you decide when to alert?''
\end{questionbox}

\begin{answerbox}
We calculate an actual probability number for each close pass --- like ``1 in 50,000 chance of collision'' or ``1 in 5,000.'' Our thresholds are:
\begin{itemize}
    \item \textbf{Safe (Green)}: Crash probability less than 1 in 100,000. No action needed.
    \item \textbf{Watch (Yellow)}: Crash probability between 1 in 100,000 and 1 in 10,000. Flag for operator review.
    \item \textbf{Critical (Red)}: Crash probability above 1 in 10,000. Automatic escape plan generated immediately.
\end{itemize}
The 2009 event had a probability of approximately 1 in 5,000 --- clearly in the Critical zone. Our system would have generated the escape plan 48 hours before impact.
\end{answerbox}

\begin{tipbox}
Use the phrase ``actual probability number'' rather than ``Pc'' when speaking. Say it in plain English. If a judge is technically sharp, you can mention that the threshold of $10^{-4}$ (1 in 10,000) is the NASA/CARA operational standard.
\end{tipbox}

\newpage
\subsection*{Q5: What is TLE data and where does it come from?}

\begin{questionbox}
``You mentioned fetching satellite data from CelesTrak. What exactly is that data? Is it real?''
\end{questionbox}

\begin{answerbox}
Every satellite's orbit is officially tracked and published as a ``Two-Line Element set'' --- basically two lines of coded text that describe the satellite's orbit. Space agencies like the US Space Force generate these from radar tracking data and publish them freely.

CelesTrak is a well-respected public database run by Dr.\ T.S. Kelso, a well-known aerospace engineer, that mirrors and organizes this official tracking data. It is the same data used by universities, research institutions, and commercial operators worldwide.

We connect to this database, download the latest tracking data for active satellites, and then use physics equations to calculate exactly where each satellite is right now, and where it will be in the next 72 hours.
\end{answerbox}

\begin{tipbox}
Show the live data table in the demo (Step 1). Say: ``This data is live. These are real satellite names and real positions right now.'' Clicking refresh in front of the judges makes it obvious it is live. That is very impressive.
\end{tipbox}

\newpage
%% ============================================================
\section{Questions About Physics and Math}
%% ============================================================

\subsection*{Q6: How do you calculate where a satellite is right now?}

\begin{questionbox}
``Exactly how does your system figure out the position of a satellite?''
\end{questionbox}

\begin{answerbox}
We use a physics model called \textbf{SGP4} --- it stands for ``Simplified General Perturbations 4'' and is the international aerospace standard for this calculation. It has been used since the 1980s.

The input is the satellite's TLE data (its official orbit description). The output is the satellite's exact X, Y, Z position and velocity in space at any moment in time. It accounts for the fact that Earth is not a perfect sphere, and that there is a tiny bit of atmosphere even in low orbit that gradually slows satellites down.

For our purposes (predicting positions up to 72 hours ahead), it is accurate to within a few hundred metres --- which is perfectly sufficient for screening.
\end{answerbox}

\begin{tipbox}
A good way to explain this: ``Think of it like a GPS. GPS satellites broadcast their position using physics equations. We use a similar, more advanced equation to calculate a satellite's position from its orbit description.''
\end{tipbox}

\vspace{1em}
\subsection*{Q7: How do you avoid checking every satellite against every other satellite (which would be too slow)?}

\begin{questionbox}
``With 27,000 satellites, checking every pair against each other would be 360 million comparisons. How do you handle that?''
\end{questionbox}

\begin{answerbox}
We use a two-step filter:

\textbf{Step 1 (Height check):} Satellites can only collide if they fly at a similar height. We instantly discard all pairs where the two satellites orbit at very different heights. This eliminates over 90\% of pairs in milliseconds.

\textbf{Step 2 (Precise check):} For the small number of pairs that pass the height filter, we calculate exactly how close the two satellites will get over the next 72 hours, and at what time.

We also have an \textbf{ML pre-filter} (explained in Q10) that quickly ranks the remaining pairs by estimated risk before running the full calculation, so the expensive physics only runs on the top suspects.
\end{answerbox}

\begin{tipbox}
Good analogy: ``It is like sorting mail. First by city (fast, eliminates most). Then by street (more detail). Then find the exact house. We sort satellites by height first, then by exact path, then by ML score.''
\end{tipbox}

\newpage
\subsection*{Q8: How do you calculate the actual probability of a crash?}

\begin{questionbox}
``You say you calculate a crash probability. How exactly is that number computed?''
\end{questionbox}

\begin{answerbox}
We know two things: how close the satellites will get, and how uncertain our position measurements are (satellite tracking has small errors of a few hundred metres).

We model these uncertainties as a ``bell curve'' spread around the predicted position. If the two satellites are predicted to pass very close together AND our measurement uncertainty is large, the chance of an actual collision is higher.

The formula combines these factors into a single probability number. It is an established aerospace standard formula, and it computes in microseconds.

For the 2009 event, the two satellites were predicted to pass within 3 metres of each other, with measurement uncertainties of about 500 metres. That gives a 1-in-5,000 crash probability --- which our system correctly flags as Critical.
\end{answerbox}

\begin{tipbox}
If a judge asks you to write the formula: $P_c \approx \frac{r^2}{2\sigma^2} e^{-d^2/(2\sigma^2)}$ where $r$ is the satellite size, $d$ is the closest distance, and $\sigma$ is the measurement uncertainty. Keep it simple and explain each variable.
\end{tipbox}

\vspace{1em}
\subsection*{Q9: Why does acting 24 hours early use 14 times less fuel?}

\begin{questionbox}
``You claim a 14$\times$ fuel saving by acting early. Explain the physics behind this.''
\end{questionbox}

\begin{answerbox}
When you push a satellite slightly in the direction it is already flying, you change its orbital height very slightly. A higher orbit moves slightly slower. Over time, this means the satellite gradually falls behind --- it drifts along its path relative to where it would have been.

Here is the key: this drift keeps building up over time. After 24 hours, the satellite has drifted 14 times further away compared to if you had applied the same push just 1 hour before the collision.

It is like steering a ship. A tiny rudder adjustment 10 km before a reef sends you safely around it. The same tiny adjustment 100 metres from the reef is not enough. Same idea with satellites --- early action means less fuel and more safety.
\end{answerbox}

\begin{tipbox}
In the demo, drag the ``how early we fire'' slider from 1 hour to 24 hours and show the gap number jumping up. This is the most visual and impressive part of the demo. Let the number speak for itself.
\end{tipbox}

\newpage
%% ============================================================
\section{Questions About the ML Model}
%% ============================================================

\subsection*{Q10: You mentioned Machine Learning. What exactly does the ML model do? Did you actually train it?}

\begin{questionbox}
``There is ML in your comparison table. What does it actually do? Did you really train a model?''
\end{questionbox}

\begin{answerbox}
Yes --- we actually trained a machine learning model. Here is exactly what it does and how:

\textbf{The problem it solves:} After our height filter, we still have hundreds of satellite pairs to check carefully. Running the full crash probability calculation on every single one is slow.

\textbf{What the model does:} Our \textbf{Random Forest model} takes cheap-to-calculate features for each pair --- things like: how close they will get, how fast they are moving toward each other, how big each satellite is, and what height they orbit at --- and produces a quick risk ranking. We then only run the expensive full crash probability calculation on the top-ranked pairs (the most suspicious ones). This makes the whole pipeline much faster without losing accuracy on real threats.

\textbf{How we trained it:} We generated our own synthetic training data. We swept through thousands of combinations of close-pass distance, relative speed, and satellite size, computed the correct crash probability for each combination using our physics formula, and trained the Random Forest on those pairs. This is the same ``surrogate model'' approach used in real aerospace conjunction screening pipelines.

\textbf{Why not use AI for the final probability?} Because crash probability must be a verifiable mathematical number. If an AI guesses wrong and a satellite gets destroyed, that is catastrophic. We use ML only to rank candidates fast --- the final probability always comes from the verified physics formula.
\end{answerbox}

\begin{tipbox}
This is a question you should answer with full confidence --- you actually built and trained this. Say: ``Yes, we trained a Random Forest on synthetic training data we generated ourselves. It pre-ranks satellite pairs by estimated risk so the expensive physics calculation only runs on the most suspicious ones. The ML ranks --- the physics proves.'' Judges will be impressed that you actually trained a model rather than just claiming AI.
\end{tipbox}

\newpage
%% ============================================================
\section{Questions About the Software}
%% ============================================================

\subsection*{Q11: How fast is your system?}

\begin{questionbox}
``Give me specific numbers. How fast does your system process satellites?''
\end{questionbox}

\begin{answerbox}
\begin{itemize}
    \item \textbf{Height filter (Step 1)}: Under 5 milliseconds for 1,000 satellites.
    \item \textbf{Precise closest-distance check}: Runs only on the small subset that passes the height filter. Fraction of a second per pair.
    \item \textbf{ML pre-ranking}: Instant --- the model scores all pairs in milliseconds.
    \item \textbf{Crash probability calculation}: Under 10 microseconds per pair. Essentially instant.
    \item \textbf{Escape plan calculation}: Under 50 microseconds per event. Essentially instant.
    \item \textbf{Full pipeline for 100 satellites}: Well under 5 seconds end-to-end.
\end{itemize}
Compare this to legacy systems that can take minutes to hours for the same catalog size.
\end{answerbox}

\begin{tipbox}
You can demonstrate this live. In the demo, click ``Refresh'' on Step 1. The data appears in 1--2 seconds (including network time). The physics calculation itself is sub-second. This is your real-time demo proof.
\end{tipbox}

\vspace{1em}
\subsection*{Q12: How do you prove the demo is actually live, not pre-recorded?}

\begin{questionbox}
``How do I know this is not just a pre-recorded video or hardcoded fake data?''
\end{questionbox}

\begin{answerbox}
There are three things you can do to verify it is live:
\begin{enumerate}
    \item \textbf{Hit Refresh in Step 1}: The satellite position coordinates update to the current UTC time. The positions change continuously because the satellites are actually moving.
    \item \textbf{Drag the sliders in Step 4}: The safety gap number updates in real time as you move the sliders. This is a live calculation, not a lookup table.
    \item \textbf{Ask us to show the terminal}: The backend server is running locally. You can see the API calls happening in real time in the terminal window.
\end{enumerate}
\end{answerbox}

\begin{tipbox}
Before the presentation, have a terminal window open with the backend logs visible. If a judge asks this question, offer to show the terminal. Live log output is very convincing proof that the system is actually calculating.
\end{tipbox}

\newpage
\subsection*{Q13: What happens when two active satellites are both trying to dodge each other at the same time?}

\begin{questionbox}
``If both satellites are active and both try to maneuver, could they accidentally maneuver into each other's escape path?''
\end{questionbox}

\begin{answerbox}
In practice, most close passes involve one active satellite and one piece of dead debris that cannot maneuver at all (like Cosmos 2251 in 2009). So usually only one satellite moves.

When both are active, real-world operators follow international coordination protocols to decide which satellite moves and in which direction.

In our current prototype, we generate an escape plan for one satellite at a time. In our Phase 3 roadmap, we plan to add multi-satellite coordination --- where the system simulates all planned maneuvers together to make sure none of them create new accidental close passes.
\end{answerbox}

\begin{tipbox}
It is totally fine to say ``our current version handles one satellite maneuver at a time. Multi-satellite coordination is on our Phase 3 roadmap.'' Judges do not expect a hackathon prototype to solve every edge case. Acknowledging limitations honestly is a sign of maturity.
\end{tipbox}

\newpage
%% ============================================================
\section{Questions About Impact and Future}
%% ============================================================

\subsection*{Q14: What would it take to deploy this for real operational use?}

\begin{questionbox}
``Could ISRO or a satellite company actually use this system as-is? What is missing?''
\end{questionbox}

\begin{answerbox}
Not as-is --- this is a research prototype. For real operational use, we would need:
\begin{itemize}
    \item \textbf{Access to better data}: Full satellite position uncertainty data (covariance matrices from Space-Track CDMs), not just public TLE files. This would make crash probability numbers more accurate.
    \item \textbf{Full catalog scaling}: GPU-accelerated processing to handle all 27,000+ satellites (not just our demo set of 100).
    \item \textbf{Security and reliability}: Real-time monitoring runs 24/7, requires backup servers, proper authentication, and uptime guarantees.
    \item \textbf{Integration}: Direct interface to satellite operators' mission control systems.
\end{itemize}
The core physics, math, ML model, and algorithms are all working correctly on real data. The jump to operational deployment is engineering scale-up --- achievable with resources.
\end{answerbox}

\begin{tipbox}
Say: ``This is a validated proof-of-concept. The hard part --- the physics, the ML, and the algorithms --- is working correctly on real historical data. The remaining steps are scale-up engineering, which is straightforward with resources.''
\end{tipbox}

\vspace{1em}
\subsection*{Q15: How does Space-Guard connect to ISRO's work?}

\begin{questionbox}
``How does this align with what ISRO is already doing in space safety?''
\end{questionbox}

\begin{answerbox}
ISRO has a dedicated space safety center called \textbf{IS4OM} (ISRO System for Safe and Sustainable Space Operations Management) at ISTRAC, Bengaluru. It is run by a dedicated directorate (DSSAM) and tracks threats to Indian satellites like Cartosat, GSAT, RISAT, and EOS.

Space-Guard is exactly the kind of automated rapid-screening and maneuver planning tool that IS4OM needs. Our system could sit upstream of IS4OM --- automatically filtering routine low-risk events and flagging only the real threats with instant escape plans, reducing the workload on ISRO operators dramatically.
\end{answerbox}

\begin{tipbox}
Focus on ISRO when talking to SIH judges --- it is the most relevant connection. Say: ``ISRO has IS4OM. Our system is designed to be exactly the kind of automated front-end tool they need.'' This shows you did your research.
\end{tipbox}

\newpage
\subsection*{Q16: What is Kessler Syndrome and how does Space-Guard help prevent it?}

\begin{questionbox}
``What is Kessler Syndrome? Is it actually a realistic threat?''
\end{questionbox}

\begin{answerbox}
Kessler Syndrome is a scenario where one collision in orbit creates so much debris that the debris causes more collisions, which create even more debris, starting a chain reaction. Eventually, an entire orbital altitude band becomes too dangerous to use.

It is named after NASA scientist Donald Kessler who described it in 1978. It is not science fiction --- scientists consider it a real risk, especially in Low Earth Orbit where satellites are most densely packed.

The 2009 Iridium--Cosmos crash was a small example: it added over 2,000 pieces of debris to an already crowded zone. If this keeps happening, eventually that altitude range becomes too dangerous for new satellites, cutting off satellite internet, GPS, weather monitoring, and more.

Space-Guard helps prevent this by catching dangerous close passes early and automatically providing the escape plan --- stopping collisions before they happen, before the debris chain reaction can begin.
\end{answerbox}

\begin{tipbox}
Say it seriously: ``Kessler Syndrome is one of the most serious long-term threats to humanity's access to space. Every prevented collision matters. Space-Guard is one piece of the solution.''
\end{tipbox}

\end{document}
